\usection{Alcance}

El presente trabajo tiene como objetivo desarrollar un lenguaje de programación visual basado en el paradigma de flujos de datos, destinado a facilitar la formación del pensamiento computacional en niños y jóvenes de entre 8 y 16 años. El proyecto abarca las etapas de análisis, diseño, implementación, validación y documentación del lenguaje.

La fase de análisis se centra en la recolección y análisis de los requerimientos que sirven de guía para el diseño del lenguaje de programación y su entorno. En esta etapa se definen las funciones básicas del lenguaje, además de las características y restricciones del diseño y arquitectura del mismo. Al final de esta fase, se obtiene como resultado el listado de requerimientos funcionales y no funcionales.

La fase de diseño del proyecto busca definir la sintaxis y la gramática del lenguaje a desarrollar. Además, se realiza el diseño de la arquitectura del entorno de ejecución, y de la interfaz visual del entorno de desarrollo, buscando su facilidad de uso para la población objetivo. Se obtiene como resultado una documentación del diseño, incluyendo diagramas, prototipos visuales de la interfaz, y la descripción de los nodos del lenguaje.

La fase de implementación tiene como objetivo desarrollar el software correspondiente al entorno de ejecución del lenguaje, y al entorno de desarrollo visual basado en flujos. El desarrollo tiene como meta cumplir con las funcionalidades descritas en las etapas anteriores. Al culminar esta etapa, se obtiene un producto mínimo viable (PMV) del lenguaje de programación. El objetivo es construir un prototipo del software que contemple las funcionalidades básicas del lenguaje, priorizando la implementación del paradigma de flujos de datos, en lugar de considerar la integración avanzada de tecnologías externas, como pudieran ser la codificación de robots o de inteligencia artificial.

En la fase de validación, se evalúa el cumplimiento de los requerimientos planteados en la primera fase, según las funcionalidades y características finales del lenguaje desarrollado. Esto se logra mediante pruebas funcionales ejecutadas por los desarrolladores, además de la validación por expertos en el área. Esta validación de expertos permitirá tener una perspectiva sobre la efectividad del lenguaje como herramienta para fomentar el desarrollo del pensamiento computacional en niños y jóvenes. Adicionalmente, se realiza una validación con usuarios de la población objetivo, con la meta de obtener retroalimentación exclusivamente sobre la usabilidad del lenguaje. Como resultado, se obtiene un reporte de los resultados del proceso para formular recomendaciones para posibles mejoras.

Por último, la fase de documentación tiene como resultado la elaboración de un manual de sistema que describa las características técnicas del lenguaje y su entorno, y un manual de usuario que brinde instrucciones de uso y guiatura general para usuarios finales.
